\subsection*{Writing answers from a pictured conversation}
A reader that returns a contiguous interval cannot combine individual words from different parts of a page. We trained a recurrent writer to choose a visible word, spell a word from its glyphs, or end its answer. Its input included the pictured passage, the current question and the preceding human exchange. A second writer learned to predict a continuous representation of the next printed word before spelling it. Neither used a text model during inference.

The direct writer matched the normalized first human answer on 1,256 of 4,413 conversation turns from reserved stories and 1,747 of 5,588 turns from eligible official development stories (Fig.~\ref{fig:dialogue-ink}). Exchanging the question and history within a story reduced these counts to 149 and 234. Replacing the passage reduced them to 492 and 675. The span reader matched 541 and 768. When we substituted a single year for its successor in a passage, the direct writer returned both the original and replacement years on 26/47 reserved pairs and 31/50 development pairs. These interventions show that some answers track the particular words on the page.

Reconstruction training improved the conditional spelling of a supplied word image. On development, after supplying the correct previous words and selecting the spelling route, the word-code writer spelled 299/599 required words from their target visual features. Predicting these features from the dialogue reduced the count to 5/599. The corresponding reserved counts were 306/573 and 7/573. One development reply, \emph{he was only a cat}, reproduced a previously unobserved answer phrase by selecting five visible words. Among the 478 reserved and 570 development turns whose normalized first answer required an absent multiword phrase, the writer produced zero and one exact answers. The direct writer produced zero in both groups.

Short answers admit alternatives. Against four human references per development turn, the direct writer obtained 46.7\% official word F1 and 38.6\% exact match on the eligible subset. The word-code writer obtained 46.2\% and 38.2\%. In 83 previously unseen-answer cases, all four references required absent multiword phrasing. Neither writer matched any of these references. Word F1 was 23.4\% for the direct writer, 25.3\% for the word-code writer and 16.5\% for the span reader. Prediction of contextually appropriate content remains the limiting step even when individual printed words can be rendered.
